%%%%%%%%%%%%%%%%%
% This is an sample CV template created using altacv.cls
% (v1.7.2, 28 August 2024) written by LianTze Lim (liantze@gmail.com).
%%%%%%%%%%%%%%%%%

\documentclass[8pt,a4paper,ragged2e,withhyper]{altacv}

\geometry{left=1.25cm,right=1.25cm,top=.5cm,bottom=1.cm,columnsep=1.2cm}

\usepackage{paracol}
\usepackage{fontawesome5}
\usepackage{hyperref}

\iftutex 
  \setmainfont{Roboto Slab}
  \setsansfont{Lato}
  \renewcommand{\familydefault}{\sfdefault}
\else
  \usepackage[rm]{roboto}
  \usepackage[defaultsans]{lato}
  \renewcommand{\familydefault}{\sfdefault}
\fi

\definecolor{SlateGrey}{HTML}{2E2E2E}
\definecolor{LightGrey}{HTML}{666666}
\definecolor{DarkPastelRed}{HTML}{8AC0D9}
\definecolor{PastelRed}{HTML}{00B0DA}
\definecolor{GoldenEarth}{HTML}{B4AD0C}

\colorlet{name}{black}
\colorlet{tagline}{PastelRed}
\colorlet{heading}{DarkPastelRed}
\colorlet{headingrule}{GoldenEarth}
\colorlet{subheading}{PastelRed}
\colorlet{accent}{PastelRed}
\colorlet{emphasis}{SlateGrey}
\colorlet{body}{LightGrey}

\definecolor{violet}{HTML}{5A2582}
\definecolor{rouge}{HTML}{F53B3B}
\definecolor{noir}{HTML}{00232C}
\definecolor{bleu}{HTML}{00B0DA}
\definecolor{rose}{HTML}{F831A0}
\definecolor{jaune}{HTML}{FFEC00}
\definecolor{vert}{HTML}{34AD6C}

\renewcommand{\namefont}{\Huge\rmfamily\bfseries}
\renewcommand{\personalinfofont}{\footnotesize}
\renewcommand{\cvsectionfont}{\LARGE\rmfamily\bfseries}
\renewcommand{\cvsubsectionfont}{\large\bfseries}

\renewcommand{\cvItemMarker}{{\small\textbullet}}
\renewcommand{\cvRatingMarker}{\faCircle}

\input{../../_shared/pubs-num.tex}
\addbibresource{../../_shared/sample.bib}

\begin{document}

\name{BOILEAU Guillaume}
\tagline{RAMS / Safety-Oriented Systems Engineer | Space Systems, Reliability, Risk Analysis \& Validation}

\photoL{2.8cm}{../../_shared/moi}

\personalinfo{%
  \email{guillaume.boileaupro@gmail.com}
  \phone{+33 6 59 94 85 84}
  \location{Alpes-Maritimes, FRANCE}
  \linkedin{guillaume-boileau-phd-891b81265}
  \researchgate{Guillaume-Boileau-2}
  \orcid{0000-0002-3576-6968}
}

\makecvheader

\vspace{-0.3cm}
\AtBeginEnvironment{itemize}{\small}

\columnratio{0.64}

\begin{paracol}{2}

\cvsection{Experience}

\cvevent{\textbf{Software Engineer – System Integration \& Hardware Interfaces}}{VEV Platform Services France}{2025 -- 2026}{Valbonne, France}

\textbf{Context:} Development and integration of software services for EV charging infrastructure within a CPMS platform connected to operational hardware systems.

\textbf{Objective:} Deliver reliable software components, robust operational data flows and maintainable service interactions with charging station systems.

\begin{itemize}
\item \textbf{System integration:} Contributed to the integration and maintenance of software services interacting with EV charging infrastructure and operational hardware.
\item \textbf{Reliability of operations:} Supported verification, monitoring and correction of operational data exchanges to ensure service continuity and consistency.
\item \textbf{Technical issue analysis:} Investigated software and hardware interaction issues in production environments and contributed to issue resolution.
\item \textbf{Data flow robustness:} Implemented and monitored FTP-based services for operational data exchange between field systems and platform services.
\item \textbf{Software development:} Developed web and backend services using TypeScript, Angular and Node.js.
\end{itemize}

\divider

\cvevent{\textbf{Senior R\&D Engineer – Space Systems, Reliability Analysis \& Validation}}{Laboratoire Lagrange, Observatoire de la Côte d'Azur, CNES}{2023 -- 2025}{Nice, France}

\textbf{Context:} Engineering and R\&D activities for the LISA ESA/NASA space mission, involving complex analysis chains, mission-level performance studies and international technical coordination.

\textbf{Objective:} Develop, structure and validate robust software and analysis workflows for complex space mission systems under strong performance and consistency constraints.

\begin{itemize}
\item \textbf{Space mission environment:} Contributed to technical activities for the LISA space mission, a large-scale ESA/NASA mission involving complex system interfaces, scientific requirements and mission-level performance constraints.
\item \textbf{Reliability and performance analysis:} Developed methods and workflows to assess the robustness, consistency and performance of signal-processing chains under uncertainty and low signal-to-noise conditions.
\item \textbf{Risk-oriented technical analysis:} Investigated model inconsistencies, unexpected behaviours and sensitivity limitations to identify potential sources of degradation and support corrective actions.
\item \textbf{Verification and validation:} Structured comparison campaigns, validation workflows and analysis scenarios to verify consistency between models, outputs and expected mission-level behaviour.
\item \textbf{Software platform development:} Developed a Python-based platform for large-scale model integration, comparison and analysis workflows. \textcolor{DarkPastelRed}{\href{https://gitlab.oca.eu/gboileau/lisa_l3_tools}{\bf GitLab:CATpyLISA}}
\item \textbf{Technical coordination:} Coordinated software development tasks, aligned contributors and supported collaborative engineering workflows across scientific and technical teams.
\item \textbf{Team management:} Supervised interns and supported engineering activities across development, analysis and validation tasks.
\end{itemize}

\divider

\cvevent{\textbf{Data Scientist – Noise, Risk \& Performance Modeling for Precision Instruments}}{Universiteit Antwerpen, Belgium}{2021 -- 2023}{Antwerp, Belgium}

\textbf{Context:} Data analysis and performance studies for precision instruments in a high-requirement technical environment.

\textbf{Objective:} Identify, characterize and mitigate sources affecting system sensitivity, reliability of measurements and system-level performance.

\begin{itemize}
\item \textbf{Noise and sensitivity analysis:} Developed robust analysis pipelines for noise source identification, characterization and mitigation.
\item \textbf{Performance degradation studies:} Investigated environmental and instrumental effects affecting measurement quality and system-level performance.
\item \textbf{Risk reduction:} Supported the identification of critical contributors to sensitivity loss and proposed analysis strategies to reduce residual uncertainties.
\item \textbf{Validation workflows:} Produced structured technical analyses requiring traceability, reproducibility and reliability of results.
\item \textbf{Technical communication:} Reported complex findings in an international research and engineering environment.
\end{itemize}

\divider

\cvevent{\textbf{Junior R\&D Engineer – Signal Analysis, Simulation \& System Investigations}}{ARTEMIS, Observatoire de la Côte d'Azur}{2018 -- 2021}{Nice, France}

\textbf{Context:} Engineering and analysis activities for gravitational-wave mission preparation, involving signal analysis, simulation and noise characterization.

\textbf{Objective:} Develop robust analysis methods for complex signals and support technical investigations in demanding system environments.

\begin{itemize}
\item \textbf{Complex signal analysis:} Developed methods for the separation of overlapping low-SNR signals in complex datasets.
\item \textbf{Simulation studies:} Built simulation approaches for complex signal environments, uncertainty studies and large-scale datasets.
\item \textbf{System diagnostics:} Investigated noise sources through cross-sensor correlation analyses in diagnostic contexts.
\item \textbf{Mission design impact:} Contributed to studies that helped identify fundamental noise and design constraints for future gravitational-wave observatories.
\item \textbf{Collaborative technical work:} Contributed to analysis methods and software tools for complex space-system applications.
\end{itemize}

\switchcolumn

\cvsection{Summary and Strengths}

\textbf{Engineer with a strong background in space systems, reliability-oriented analysis, validation and technical risk investigation}, developed through major international scientific and engineering collaborations including \textbf{LISA}, \textbf{LIGO--Virgo--KAGRA} and \textbf{Einstein Telescope}.

Experienced in \textbf{complex system analysis}, \textbf{performance verification}, \textbf{model validation}, \textbf{technical diagnostics}, \textbf{noise and uncertainty analysis}, and \textbf{robust software workflows}. Used to working in demanding environments requiring \textbf{rigor}, \textbf{traceability}, \textbf{structured execution}, \textbf{technical reporting} and \textbf{cross-functional coordination}.

Relevant background for \textbf{RAMS and Safety activities}: reliability analysis, sensitivity studies, risk-oriented investigations, identification of critical contributors, validation against requirements, residual uncertainty assessment and recommendations for system improvement.

Additional hands-on experience in \textbf{software development and system interfaces} for EV charging infrastructure using TypeScript, Angular and Node.js, including operational data exchange services and production issue analysis.

\cvsection{RAMS / Safety Relevance}

\begin{itemize}
  \item Experience with \textbf{complex space mission environments}, especially LISA, involving strict performance constraints and system-level consistency requirements.
  \item Strong background in \textbf{reliability-oriented analysis}, including robustness studies, sensitivity limitations, uncertainty propagation and performance degradation analysis.
  \item Ability to identify \textbf{technical risks, critical contributors and residual uncertainties} through data-driven investigations and model comparisons.
  \item Experience in \textbf{verification and validation workflows}, including structured campaigns, consistency checks and technical reporting.
  \item Comfortable interfacing with multidisciplinary teams in \textbf{English-speaking international collaborations}.
\end{itemize}

\cvsection{Team Management}

\begin{itemize}
  \item Supervision of \textbf{3 Master's thesis interns (M2)} during engineering and development phases.
  \item Coordination of technical activities involving \textbf{research engineers and technicians} on a \textbf{data processing pipeline} for the \textbf{LISA space mission}.
  \item Experience in organizing tasks, supporting contributors and maintaining alignment across collaborative technical developments in demanding environments.
\end{itemize}

\cvsection{Languages}

\cvskill{English}{4}
\divider

\cvskill{Spanish}{2}
\divider

\medskip

\cvsection{Education}

\cvevent{PhD in Physics}{Université Côte d'Azur}{2018 -- 2021}{Nice, France}
\divider

\cvevent{Selective Graduate Program in Fundamental Physics (Magistère)}{Université Paris-Saclay}{2016 -- 2018}{Orsay, France}
\divider

\cvevent{MSc in Astronomy and Astrophysics}{Observatoire de Paris / Université Paris-Saclay}{2017 -- 2018}{Paris, France}
\divider

\cvevent{BSc in Fundamental Physics}{Université Paris-Saclay}{2015 -- 2016}{Orsay, France}

\cvsection{Professional Training}

\cvevent{Supervised Machine Learning (Regression \& Classification)}{DeepLearning.AI – Andrew Ng}{2020}{Coursera}

\end{paracol}

\newpage

\cvsection{Projects}

\cvevent{Co-author and full member of the LISA Consortium}{ESA/NASA Space Mission -- Large-scale International Collaboration}{2018 -- now}{}

Contribution to software, analysis and validation workflows for the LISA space mission within an international engineering and scientific collaboration involving ESA, NASA and multiple research institutes.

\textbf{Role:} Contribution to software workflows, technical coordination, complex signal-processing activities, model validation and collaborative engineering tasks.

\textbf{RAMS relevance:} Work on system-level performance, signal sensitivity, robustness of analysis chains, uncertainty assessment and consistency verification in a mission context with strong reliability and validation constraints.

\textbf{Program scale:} Major international space mission with an estimated budget of approximately 2 billion EUR over 10 years.

\divider

\cvevent{Co-author and full member of Einstein Telescope and LIGO--Virgo--KAGRA Collaborations}{International Collaborations}{2021 -- now}{}

Contribution to collaborative developments in data analysis, signal characterization, software methods and technical investigations within the LIGO--Virgo--KAGRA and Einstein Telescope collaborations.

\textbf{Role:} Development of analysis tools, contribution to technical activities and support to complex software workflows in high-standard collaborative environments.

\textbf{System impact:} Contributed to studies on environmental and instrumental noise sources, helping identify fundamental design constraints and performance limitations for next-generation gravitational-wave observatories.

\textbf{Environment:} Large-scale international collaborations involving strong technical requirements, structured methods, technical review and coordinated contributions.

\cvsection{Strengths}

\vspace{-.3cm}

\cvsubsection{Core skills}

{\small
\cvtag{RAMS-Oriented Analysis}
\cvtag{Reliability Analysis}
\cvtag{Safety-Oriented Engineering}
\cvtag{Risk Analysis}
\cvtag{Residual Risk Assessment}
\cvtag{Criticality Analysis}
\cvtag{Failure Scenario Analysis}
\cvtag{System-Level Analysis}
\cvtag{Requirements Verification}
\cvtag{Verification \& Validation}
\cvtag{System Validation}
\cvtag{Integration Testing}
\cvtag{Performance Verification}
\cvtag{Technical Investigation}
\cvtag{Anomaly Investigation}
\cvtag{Root Cause Analysis}
\cvtag{Technical Diagnostics}
\cvtag{Issue Resolution}
\cvtag{Traceability}
\cvtag{Structured Execution}
\cvtag{Technical Reporting}
\cvtag{Complex Systems}
\cvtag{Space Systems}
\cvtag{Signal Processing}
\cvtag{Noise Analysis}
\cvtag{Uncertainty Analysis}
\cvtag{Software Integration}
\cvtag{Technical Coordination}
}

\divider\smallskip
\vspace{-.3cm}

\cvsubsection{Technical skills}

{\small
\cvtag{Python}
\cvtag{Git}
\cvtag{Docker}
\cvtag{Linux}
\cvtag{SQL}
\cvtag{TypeScript}
\cvtag{Angular}
\cvtag{Node.js}
\cvtag{C}
\cvtag{Matlab}
\cvtag{Pandas}
\cvtag{REST API}
\cvtag{Bash scripting}
\cvtag{FTP}
\cvtag{Test Automation Scripting}
\cvtag{CI/CD}
\cvtag{Issue Tracking}
\cvtag{Technical Reporting}
\cvtag{Condor}
\cvtag{Slurm}
\cvtag{PyTorch}
\cvtag{MongoDB}
\cvtag{Mathematica}
}

\divider\smallskip
\vspace{-.3cm}

\cvsubsection{Domain knowledge}

{\small
\cvtag{Satellite Mission Context}
\cvtag{ESA/NASA Collaboration}
\cvtag{Mission-Level Performance}
\cvtag{Scientific Payload Analysis}
\cvtag{Precision Instrumentation}
\cvtag{Environmental Noise}
\cvtag{Space Mission Data Processing}
\cvtag{Large-Scale International Projects}
\cvtag{Technical English}
}

\divider\smallskip
\vspace{-.3cm}

{\small\LaTeXraggedright
\cvsubsection{Soft skills}

{\small
\cvtag{Autonomy}
\cvtag{Analytical Thinking}
\cvtag{Execution Rigor}
\cvtag{Problem Solving}
\cvtag{Adaptability}
\cvtag{Responsibility}
\cvtag{Attention to Detail}
\cvtag{Collaboration}
\cvtag{Stakeholder Communication}
\cvtag{Structured Communication}
\cvtag{Team Coordination}
\cvtag{Customer Awareness}
\cvtag{Technical Interface}
\par}}

\cvtag{\textit{Applications: RAMS support, safety-oriented analysis, space systems, risk investigation, verification and validation workflows}}

\end{document}